\begin{definitionbox}{Definition}
\begin{definition}[Embedding of $\mathbb{Q}$ into $\mathbb{R}$]
\label{def:real-chapter-embedding-q-into-r}
Define $\iota:\mathbb{Q}\to\mathbb{R}$ by
\[
\iota(q):=q^{\ast}.
\]
\end{definition}
\end{definitionbox}

\begin{remark*}[Standard quantified statement]
\[
\iota:\mathbb{Q}\to\mathbb{R}
\quad\text{and}\quad
\forall q\in\mathbb{Q},\ \iota(q)=q^{\ast}.
\]
\end{remark*}

\begin{remark*}[Interpretation]
The canonical embedding sends each rational number to its corresponding rational cut.
\end{remark*}

\begin{dependencies}
\begin{itemize}
  \item \hyperref[def:rationals]{Rational Numbers}
  \item \hyperref[def:reals]{The Real Numbers}
  \item \hyperref[def:rational-cut]{Rational Cut}
\end{itemize}
\end{dependencies}

\begin{theorembox}{Theorem}
\begin{theorem}[Embedding Preserves Zero and One]
\label{thm:embedding-q-into-r-preserves-zero-one}
The embedding satisfies
\[
\iota(0_{\mathbb{Q}})=0_{\mathbb{R}}
\qquad\text{and}\qquad
\iota(1_{\mathbb{Q}})=1_{\mathbb{R}}.
\]
\hyperref[prf:embedding-q-into-r-preserves-zero-one]{\textit{Go to proof.}}
\end{theorem}
\end{theorembox}

\begin{remark*}[Standard quantified statement]
\[
\iota(0_{\mathbb{Q}})=0_{\mathbb{R}}
\land
\iota(1_{\mathbb{Q}})=1_{\mathbb{R}}.
\]
\end{remark*}

\begin{remark*}[Interpretation]
The rational constants zero and one become the designated real constants under the embedding.
\end{remark*}

\begin{dependencies}
\begin{itemize}
  \item \hyperref[def:real-chapter-embedding-q-into-r]{Embedding of $\mathbb{Q}$ into $\mathbb{R}$}
  \item \hyperref[def:zero-in-q]{Zero in $\mathbb{Q}$}
  \item \hyperref[def:one-in-q]{One in $\mathbb{Q}$}
  \item \hyperref[def:zero-in-r]{Zero in $\mathbb{R}$}
  \item \hyperref[def:one-in-r]{One in $\mathbb{R}$}
\end{itemize}
\end{dependencies}

\begin{theorembox}{Theorem}
\begin{theorem}[Embedding $\mathbb{Q}$ into $\mathbb{R}$ Is Injective]
\label{thm:embedding-q-into-r-injective}
For all $p,q\in\mathbb{Q}$,
\[
\iota(p)=\iota(q) \Longrightarrow p=q.
\]
\hyperref[prf:embedding-q-into-r-injective]{\textit{Go to proof.}}
\end{theorem}
\end{theorembox}

\begin{remark*}[Standard quantified statement]
\[
\forall p,q\in\mathbb{Q},\qquad
\iota(p)=\iota(q)\Longrightarrow p=q.
\]
\end{remark*}

\begin{remark*}[Predicate reading]
\[
\operatorname{Injective}(\iota)
\]
means that $\iota(p)=\iota(q)$ implies $p=q$ for rational inputs $p$ and $q$.
\end{remark*}

\begin{remark*}[Interpretation]
Distinct rational numbers determine distinct rational cuts, so the embedding loses no rational information.
\end{remark*}

\begin{dependencies}
\begin{itemize}
  \item \hyperref[def:real-chapter-embedding-q-into-r]{Embedding of $\mathbb{Q}$ into $\mathbb{R}$}
  \item \hyperref[def:rational-cut]{Rational Cut}
  \item \hyperref[thm:q-totally-ordered-set]{$\mathbb{Q}$ Is a Totally Ordered Set}
\end{itemize}
\end{dependencies}

\begin{theorembox}{Theorem}
\begin{theorem}[Embedding Preserves Addition and Multiplication]
\label{thm:embedding-q-into-r-preserves-arithmetic}
For all $p,q\in\mathbb{Q}$,
\[
\iota(p+q)=\iota(p)+\iota(q)
\qquad\text{and}\qquad
\iota(p\cdot q)=\iota(p)\cdot\iota(q).
\]
\hyperref[prf:embedding-q-into-r-preserves-arithmetic]{\textit{Go to proof.}}
\end{theorem}
\end{theorembox}

\begin{remark*}[Standard quantified statement]
\[
\forall p,q\in\mathbb{Q},\qquad
\iota(p+q)=\iota(p)+\iota(q)
\land
\iota(p\cdot q)=\iota(p)\cdot\iota(q).
\]
\end{remark*}

\begin{remark*}[Predicate reading]
\[
\operatorname{Homomorphism}(\iota,\mathbb{Q},\mathbb{R},+)
\land
\operatorname{Homomorphism}(\iota,\mathbb{Q},\mathbb{R},\cdot)
\]
means that $\iota$ preserves rational addition and multiplication as real addition and multiplication.
\end{remark*}

\begin{remark*}[Interpretation]
Addition and multiplication of rational numbers agree with addition and multiplication after passing to rational cuts.
\end{remark*}

\begin{dependencies}
\begin{itemize}
  \item \hyperref[def:real-chapter-embedding-q-into-r]{Embedding of $\mathbb{Q}$ into $\mathbb{R}$}
  \item \hyperref[def:addition-on-r]{Addition on $\mathbb{R}$}
  \item \hyperref[def:multiplication-on-r]{Multiplication on $\mathbb{R}$}
  \item \hyperref[def:addition-on-q-classes]{Addition on $\mathbb{Q}$}
  \item \hyperref[def:multiplication-on-q-classes]{Multiplication on $\mathbb{Q}$}
\end{itemize}
\end{dependencies}

\begin{theorembox}{Theorem}
\begin{theorem}[Embedding Preserves Order]
\label{thm:embedding-q-into-r-preserves-order}
For all $p,q\in\mathbb{Q}$,
\[
p\le q \Longleftrightarrow \iota(p)\le\iota(q).
\]
\hyperref[prf:embedding-q-into-r-preserves-order]{\textit{Go to proof.}}
\end{theorem}
\end{theorembox}

\begin{remark*}[Standard quantified statement]
\[
\forall p,q\in\mathbb{Q},\qquad
p\le q \Longleftrightarrow \iota(p)\le\iota(q).
\]
\end{remark*}

\begin{remark*}[Predicate reading]
\[
\operatorname{OrderEmbedding}(\iota,\mathbb{Q},\mathbb{R},\le)
\]
means that $\iota$ preserves and reflects the nonstrict order.
\end{remark*}

\begin{remark*}[Interpretation]
The order on rational cuts extends the original rational order exactly.
\end{remark*}

\begin{dependencies}
\begin{itemize}
  \item \hyperref[def:real-chapter-embedding-q-into-r]{Embedding of $\mathbb{Q}$ into $\mathbb{R}$}
  \item \hyperref[def:order-on-q]{Order on $\mathbb{Q}$}
  \item \hyperref[def:order-on-r]{Order on $\mathbb{R}$}
  \item \hyperref[def:rational-cut]{Rational Cut}
\end{itemize}
\end{dependencies}

\begin{theorembox}{Theorem}
\begin{theorem}[$\mathbb{Q}$ Embeds into $\mathbb{R}$ as an Ordered Field]
\label{thm:q-embeds-into-r}
The map $\iota$ is an injective order-preserving field homomorphism.
\hyperref[prf:q-embeds-into-r]{\textit{Go to proof.}}
\end{theorem}
\end{theorembox}

\begin{remark*}[Standard quantified statement]
\[
\operatorname{Injective}(\iota)
\land
\operatorname{Homomorphism}(\iota,\mathbb{Q},\mathbb{R},+)
\land
\operatorname{Homomorphism}(\iota,\mathbb{Q},\mathbb{R},\cdot)
\land
\operatorname{OrderEmbedding}(\iota,\mathbb{Q},\mathbb{R},\le).
\]
\end{remark*}

\begin{remark*}[Interpretation]
The rational numbers sit inside the real numbers as an ordered subfield through the rational-cut embedding.
\end{remark*}

\begin{dependencies}
\begin{itemize}
  \item \hyperref[def:real-chapter-embedding-q-into-r]{Embedding of $\mathbb{Q}$ into $\mathbb{R}$}
  \item \hyperref[thm:embedding-q-into-r-preserves-zero-one]{Embedding Preserves Zero and One}
  \item \hyperref[thm:embedding-q-into-r-injective]{Embedding $\mathbb{Q}$ into $\mathbb{R}$ Is Injective}
  \item \hyperref[thm:embedding-q-into-r-preserves-arithmetic]{Embedding Preserves Addition and Multiplication}
  \item \hyperref[thm:embedding-q-into-r-preserves-order]{Embedding Preserves Order}
\end{itemize}
\end{dependencies}

\begin{theorembox}{Theorem}
\begin{theorem}[Density of $\mathbb{Q}$ in $\mathbb{R}$]
\label{thm:density-of-q-in-r}
For all $\alpha,\beta\in\mathbb{R}$,
\[
\alpha<\beta \Longrightarrow \exists q\in\mathbb{Q}\,(\alpha<\iota(q)<\beta).
\]
\hyperref[prf:density-of-q-in-r]{\textit{Go to proof.}}
\end{theorem}
\end{theorembox}

\begin{remark*}[Standard quantified statement]
\[
\forall \alpha,\beta\in\mathbb{R},\qquad
\alpha<\beta\Longrightarrow
\exists q\in\mathbb{Q}\,(\alpha<\iota(q)\land \iota(q)<\beta).
\]
\end{remark*}

\begin{remark*}[Interpretation]
Every nonempty open interval of real numbers contains an embedded rational number.
\end{remark*}

\begin{dependencies}
\begin{itemize}
  \item \hyperref[def:real-chapter-embedding-q-into-r]{Embedding of $\mathbb{Q}$ into $\mathbb{R}$}
  \item \hyperref[def:strict-order-on-r]{Strict Order on $\mathbb{R}$}
  \item \hyperref[thm:density-of-q]{Density of $\mathbb{Q}$ in Itself}
  \item \hyperref[def:dedekind-cut]{Dedekind Cut}
\end{itemize}
\end{dependencies}

\begin{theorembox}{Theorem}
\begin{theorem}[Archimedean Property of $\mathbb{R}$]
\label{thm:archimedean-property-of-r}
For $\alpha,\beta\in\mathbb{R}$,
\[
0_{\mathbb{R}}<\alpha
\Longrightarrow
\exists n\in\mathbb{N}\,(\beta<\iota(n)\cdot\alpha).
\]
\hyperref[prf:archimedean-property-of-r]{\textit{Go to proof.}}
\end{theorem}
\end{theorembox}

\begin{remark*}[Standard quantified statement]
\[
\forall \alpha,\beta\in\mathbb{R},\qquad
0_{\mathbb{R}}<\alpha
\Longrightarrow
\exists n\in\mathbb{N}\,(\beta<\iota(n)\cdot\alpha).
\]
\end{remark*}

\begin{remark*}[Interpretation]
Positive real multiples of a sufficiently large natural number eventually exceed any fixed real number.
\end{remark*}

\begin{dependencies}
\begin{itemize}
  \item \hyperref[thm:lub-property-r]{Least-Upper-Bound Property}
\end{itemize}
\end{dependencies}
